\section{An adaptive counterexample}

The following lower bound is for the algorithm specified above and the ideal
independent fair-bit model. It does not concern the distinct algorithms
analyzed by Brodal. The adversary sees the current tree before choosing each
next key, and all choices are made before the next insertion's random bits.

\begin{theorem}[Linear adaptive depth and repair]\label{thm:adaptive-lower}
There is an adaptive sequence of distinct rational keys and a key $x$ in the
initial prefix such that, after $n$ successful insertions,
$\mathbb E[\operatorname{depth}_n(x)]=\Omega(n)$. Along an infinite subsequence
of insertions, the expected number of counted repair visits is also $\Omega(n)$.
The expectations include all random bits, including those in the initial prefix.
\end{theorem}

\begin{proof}
Throughout, $x$ will be the largest inserted key, so its root-to-node path
consists entirely of right edges. Let $n$ include the next key, and write
$k=\lceil\log_2 n\rceil$. Assume for the moment that $k\geq2$ and that the
present depth $d$ of $x$ satisfies $d>k+2$. We specify two types of insertion.
All required gaps exist because the rational order is dense.
For definiteness, choose the midpoint of each bounded gap and one less than
the upper endpoint for an unbounded-below gap. The attack depends only on the
current tree, including its keys; the fixed schedule below is determined by
the number of nodes, which can be counted from that tree.

For an \emph{attack}, let $g$ be the node at depth $k$ on the path to $x$,
and insert a fresh key in $g$'s left subtree, smaller than every key already
in that subtree. If the subtree is empty, insert at its root gap. The insertion
path is $k$ right edges followed by one or more left edges. There are three
cases. With probability $1/2$, $T=k$: a left child of $g$ is promoted once,
so $g$, its right subtree, and $x$ move down one level. This includes the
case where the new leaf is the left child of $g$. With probability $1/4$,
$T=k-1$: the selected two-edge path is right--left, so the double rotation
leaves $g$ at its old depth and leaves $d$ unchanged. With probability $1/4$,
$T\leq k-2$: the selected path is right--right and the single rotation moves
$x$ up one level. Thus an attack changes $d$ by $+1,0,-1$ with respective
probabilities $1/2,1/4,1/4$.

For a \emph{probe}, insert a fresh key in the gap immediately below $x$,
choosing it within $x$'s left subtree. The insertion path reaches $x$, and
its insertion depth is at least $d+1$. Since $d>k+2$, every target $T\leq k$
selects a right--right path strictly above $x$. Exactly one rotation occurs,
moving $x$ up one level. By the exact visit formula, the probe has
\[ V=(d_{\mathrm{ins}}-T-1)_+\geq d-k. \]
Neither kind of key exceeds $x$, which remains the maximum.

Fix $m$ in advance. To initialize, insert $m$ distinct keys in decreasing
order, starting with $0$ and then choosing one less than the current minimum.
Afterward, the current maximum is $x$. At each subsequent step, count the
nodes to determine the next size and its position in a fixed schedule of nine
attacks followed by one probe. If the current depth of $x$ is at most $k+2$,
insert one less than the current minimum; otherwise use the scheduled attack
or probe. This is a deterministic function of the current tree and fixed $m$:
it does not remember whether the condition failed previously, inspect priorities
or random bits, or know future outcomes. On the event that
each sampled target in this prefix is zero, every new minimum rotates to the
root, leaving a right chain with $d_m(x)=m-1$. This event has strictly positive
probability for every fixed $m$.

It remains to choose a fixed $m$ so that the subsequent condition holds forever
with positive probability. Couple each future attack to an independent random
variable taking values $+1,0,-1$ with probabilities $1/2,1/4,1/4$; assign
$-1$ to each future probe. The coupling, used only in the analysis and not
available to the adversary, can read the first two latent fair bits
for an attack: zero gives $+1$, one then zero gives $0$, and two ones give
$-1$. Latent bits are independent even when a sampler stops before reading
them. Let $Z_s$ be the sum of the first $s$ such scheduled increments, with
$Z_0=0$. The sums of successive ten-increment blocks are independent and have
mean $9/4-1=5/4$. The strong law of large numbers, and the bounded remainder
within each block, imply
\[ \frac{Z_s}{s}\longrightarrow\frac18 \quad\text{almost surely}. \]
Consequently
\[ M=\sup_{s\geq1}\bigl(\lceil\log_2(s+1)\rceil-Z_{s-1}\bigr) \]
is finite almost surely. Choose a finite integer $B$ with $\Prb(M\leq B)>0$,
then choose a fixed $m\geq4$ large enough that
$m-1-\lceil\log_2m\rceil\geq B+3$. The initial-prefix event and $M\leq B$
are independent and their intersection $E$ has positive probability.

For step $s$ after the prefix, ceiling subadditivity and $m+s\leq m(s+1)$ give
\[\lceil\log_2(m+s)\rceil\leq
  \lceil\log_2m\rceil+\lceil\log_2(s+1)\rceil.\]
Inductively on $E$, before that step the depth is $m-1+Z_{s-1}$, so its
excess over the current $k$ is at least $3$. Hence the stated attack and probe
transitions apply at every step, and on $E$,
$d_{m+s}(x)=m-1+Z_s$ and $d_{m+s}(x)/(m+s)\to1/8$ almost surely.
Fatou's lemma yields
\[ \liminf_{n\to\infty}\frac{\mathbb E[d_n(x)]}{n}
   \geq \frac{\Prb(E)}8>0.\]
On the probe subsequence $n=m+10j$, the bound $V\geq d-k$, applied to the
depth just before the probe, and the same argument give
\[ \liminf_{j\to\infty}\frac{\mathbb E[V_{m+10j}]}{m+10j}
   \geq \frac{\Prb(E)}8>0.\]
Thus both stated lower bounds are unconditional.
\end{proof}

\begin{corollary}[Tight adaptive worst case]\label{cor:adaptive-tight}
Under the model of Theorem~\ref{thm:adaptive-lower}, the largest expected
node depth attainable by an adaptive adversary is $\Theta(n)$. The expected
repair visits in that theorem are $\Theta(n)$ along its probe subsequence.
\end{corollary}
\begin{proof}
Every node of every $n$-node binary tree has depth at most $n-1$, and the
counted visits for an insertion are at most its insertion depth. These
deterministic upper bounds, together with Theorem~\ref{thm:adaptive-lower},
give the claimed matching bounds.
\end{proof}

The proof uses an infinite dense key order. The Python implementation compares
keys at runtime, although its annotations say \texttt{int}; the rational model
must therefore be distinguished from the annotated interface. For any fixed
finite horizon, its finitely many possible adaptive histories use finitely many
rational keys; an order-preserving rescaling embeds those keys in the integers.
This yields the same finite-horizon worst-case bounds over integer inputs, with
the spacing allowed to depend on the horizon. Comparisons here have unit cost;
the integer spacing may require more than $O(\log n)$ bits per key, so the
embedding does not establish the result for a fixed-width integer universe or
account for arbitrary-precision arithmetic costs. The counterexample does not
prove a linear lower bound for a fixed, nonadaptive insertion sequence.
